\documentclass[a4paper,14pt]{extarticle}

% Кодировка и язык
\usepackage[T2A]{fontenc}
\usepackage[utf8]{inputenc}
\usepackage[russian]{babel}

% Шрифт Times New Roman
\usepackage{tempora}

% Поля: левое 3 см, правое 1 см, верхнее 2 см, нижнее 2 см
\usepackage[left=3cm, right=1cm, top=2cm, bottom=2cm, footskip=1.25cm, headsep=0.5cm]{geometry}

% Межстрочный интервал 1.5
\usepackage{setspace}
\onehalfspacing

% Отступ первой строки 1.25 см
\usepackage{indentfirst}
\setlength{\parindent}{1.25cm}

% Автоматические переносы
\usepackage{hyphenat}
\sloppy

% Нумерация страниц внизу по центру
\usepackage{fancyhdr}
\pagestyle{fancy}
\fancyhf{}
\fancyfoot[C]{\thepage}
\renewcommand{\headrulewidth}{0pt}
\renewcommand{\footrulewidth}{0pt}
\setlength{\headheight}{0pt}

% Таблицы
\usepackage{array}
\usepackage{longtable}
\usepackage{tabularx}
\usepackage{booktabs}

% Рисунки
\usepackage{graphicx}
\usepackage{float}

% Для кода — Courier New
\usepackage{listings}
\lstset{
  basicstyle=\small\ttfamily,
  breaklines=true,
  frame=none,
  numbers=none,
  tabsize=4,
  xleftmargin=1.25cm,
  showstringspaces=false,
  inputencoding=utf8,
  extendedchars=true,
  literate={а}{{\cyra}}1 {б}{{\cyrb}}1 {в}{{\cyrv}}1 {г}{{\cyrg}}1 {д}{{\cyrd}}1 {е}{{\cyre}}1 {ж}{{\cyrzh}}1 {з}{{\cyrz}}1 {и}{{\cyri}}1 {й}{{\cyrishrt}}1 {к}{{\cyrk}}1 {л}{{\cyrl}}1 {м}{{\cyrm}}1 {н}{{\cyrn}}1 {о}{{\cyro}}1 {п}{{\cyrp}}1 {р}{{\cyrr}}1 {с}{{\cyrs}}1 {т}{{\cyrt}}1 {у}{{\cyru}}1 {ф}{{\cyrf}}1 {х}{{\cyrh}}1 {ц}{{\cyrc}}1 {ч}{{\cyrch}}1 {ш}{{\cyrsh}}1 {щ}{{\cyrshch}}1 {ъ}{{\cyrhrdsn}}1 {ы}{{\cyrery}}1 {ь}{{\cyrsftsn}}1 {э}{{\cyrerev}}1 {ю}{{\cyryu}}1 {я}{{\cyrya}}1
}

% TikZ для диаграмм
\usepackage{tikz}
\usetikzlibrary{shapes.geometric, arrows.meta, positioning}

% Подписи
\usepackage[labelsep=endash]{caption}
\captionsetup[figure]{justification=centering, font={normalsize}, skip=5pt, position=below}
\captionsetup[table]{justification=raggedleft, singlelinecheck=false, font={normalsize}, skip=5pt, position=above}

% Гиперссылки (без цветных рамок)
\usepackage[hidelinks]{hyperref}

% Содержание
\usepackage{tocloft}
\renewcommand{\cftsecleader}{\cftdotfill{\cftdotsep}}
\renewcommand{\cftsubsecleader}{\cftdotfill{\cftdotsep}}
\renewcommand{\cftsubsubsecleader}{\cftdotfill{\cftdotsep}}

% Заголовки разделов
\usepackage{titlesec}

% Структурные элементы: прописные, по центру, полужирный, 16 пт, с новой страницы
\newcommand{\structsection}[1]{%
  \clearpage
  \phantomsection
  \begin{center}
    \fontsize{16pt}{20pt}\selectfont\bfseries\MakeUppercase{#1}
  \end{center}
  \addcontentsline{toc}{section}{#1}
  \vspace{1em}
}

% Нумерованные разделы
\titleformat{\section}[block]
  {\fontsize{16pt}{20pt}\selectfont\bfseries}
  {\thesection}{0.5em}{}
\titleformat{\subsection}[block]
  {\fontsize{14pt}{18pt}\selectfont\bfseries}
  {\thesubsection}{0.5em}{}
\titleformat{\subsubsection}[block]
  {\fontsize{14pt}{18pt}\selectfont\bfseries}
  {\thesubsubsection}{0.5em}{}

\titlespacing*{\section}{1.25cm}{1em}{0.5em}
\titlespacing*{\subsection}{1.25cm}{1em}{0.5em}
\titlespacing*{\subsubsection}{1.25cm}{1em}{0.5em}

% Без точки после номера раздела
\renewcommand{\thesection}{\arabic{section}}
\renewcommand{\thesubsection}{\thesection.\arabic{subsection}}
\renewcommand{\thesubsubsection}{\thesubsection.\arabic{subsubsection}}

% Список источников
\usepackage{enumitem}

\begin{document}

% ============================================================
% ТИТУЛЬНЫЙ ЛИСТ
% ============================================================
\thispagestyle{empty}
\begin{center}

\fontsize{14pt}{18pt}\selectfont

МИНИСТЕРСТВО НАУКИ И ВЫСШЕГО ОБРАЗОВАНИЯ\\РОССИЙСКОЙ ФЕДЕРАЦИИ

\vspace{0.5cm}

ФГБОУ ВО <<ДАГЕСТАНСКИЙ ГОСУДАРСТВЕННЫЙ\\ТЕХНИЧЕСКИЙ УНИВЕРСИТЕТ>>

\vspace{0.5cm}

Факультет информатики и вычислительной техники\\
Кафедра информатики

\vspace{2cm}

\fontsize{16pt}{20pt}\selectfont\bfseries
ОТЧЁТ\\
о научно-исследовательской работе

\vspace{1cm}

\fontsize{14pt}{18pt}\selectfont\normalfont
по теме:\\
<<Исследование методов и технологий проектирования RESTful\\
веб-приложений для управления задачами\\
на основе многослойной архитектуры>>

\end{center}

\vspace{2cm}

\begin{flushright}
\fontsize{14pt}{18pt}\selectfont
Выполнил: студент группы {[\_\_\_]}\\
Берсиров~С.\\[1cm]
Руководитель: {[\_\_\_]}
\end{flushright}

\vfill

\begin{center}
\fontsize{14pt}{18pt}\selectfont
Махачкала 2026
\end{center}

\newpage

% ============================================================
% РЕФЕРАТ
% ============================================================
\structsection{Реферат}

Отчёт о НИР содержит: 25~с., 3~рис., 9~табл., 15~источников.

\textbf{Ключевые слова:} REST API, GO, МНОГОСЛОЙНАЯ АРХИТЕКТУРА, ВЕБ-ПРИЛОЖЕНИЕ, УПРАВЛЕНИЕ ЗАДАЧАМИ, POSTGRESQL, DOCKER, CRUD, SWAGGER, OPENAPI.

Объект исследования~--- веб-приложения для управления задачами (Todo-приложения).

Предмет исследования~--- методы и технологии проектирования RESTful веб-приложений с применением многослойной архитектуры на языке программирования Go.

Цель работы~--- исследование существующих методов и технологий для проектирования и разработки RESTful веб-приложения управления задачами с применением многослойной архитектуры, а также подготовка входных данных для практической реализации.

В результате исследования проведён анализ предметной области управления задачами, выполнен обзор существующих аналогов, обоснован выбор технологий и архитектурных решений, сформулировано техническое задание, спроектирована структура базы данных и REST~API, определены тестовые сценарии для проверки стабильной работы системы.

% ============================================================
% СОДЕРЖАНИЕ
% ============================================================
\clearpage
\phantomsection
\begin{center}
\fontsize{16pt}{20pt}\selectfont\bfseries СОДЕРЖАНИЕ
\end{center}
\vspace{1em}
\renewcommand{\contentsname}{}
\tableofcontents

% ============================================================
% ТЕРМИНЫ И ОПРЕДЕЛЕНИЯ
% ============================================================
\structsection{Термины и определения}

В настоящем отчёте применяются следующие термины с соответствующими определениями.

\textbf{REST} (Representational State Transfer)~--- архитектурный стиль взаимодействия компонентов распределённого приложения в сети, основанный на принципах использования стандартных HTTP-методов для выполнения операций над ресурсами.

\textbf{API} (Application Programming Interface)~--- программный интерфейс приложения, представляющий собой набор определённых правил и спецификаций, посредством которых одна программа может взаимодействовать с другой.

\textbf{CRUD} (Create, Read, Update, Delete)~--- четыре базовые операции, используемые при работе с хранилищами данных: создание, чтение, обновление и удаление записей.

\textbf{Middleware}~--- промежуточное программное обеспечение, представляющее собой обработчик, включённый в цепочку обработки HTTP-запросов, выполняющий вспомогательные функции до или после основного обработчика.

\textbf{Docker}~--- платформа для разработки, доставки и запуска приложений в изолированных контейнерах, обеспечивающая воспроизводимость среды выполнения.

\textbf{Graceful Shutdown}~--- механизм корректного завершения работы серверного приложения, при котором прекращается приём новых запросов, но завершаются все текущие обработчики.

\textbf{Optimistic Locking} (оптимистичная блокировка)~--- стратегия управления конкурентным доступом к данным, при которой при обновлении записи проверяется, что её версия не изменилась с момента чтения.

\textbf{UUID} (Universally Unique Identifier)~--- универсальный уникальный идентификатор, представляющий собой 128-битное число, используемое для идентификации объектов в распределённых системах.

\textbf{JSON} (JavaScript Object Notation)~--- текстовый формат обмена данными, основанный на подмножестве языка JavaScript, широко используемый в REST~API для передачи структурированных данных.

\textbf{Swagger/OpenAPI}~--- спецификация для описания, создания, использования и визуализации RESTful веб-сервисов, позволяющая автоматически генерировать интерактивную документацию.

\textbf{PostgreSQL}~--- свободная объектно-реляционная система управления базами данных с открытым исходным кодом, обеспечивающая надёжность, целостность данных и поддержку расширенных SQL-запросов.

\textbf{Go} (Golang)~--- компилируемый многопоточный язык программирования, разработанный компанией Google, ориентированный на создание высокопроизводительных сетевых сервисов и серверных приложений.

\textbf{Clean Architecture} (чистая архитектура)~--- архитектурный подход, предложенный Робертом Мартином, основанный на разделении системы на слои с направлением зависимостей от внешних слоёв к внутренним, обеспечивающий независимость бизнес-логики от деталей реализации.

% ============================================================
% ОБОЗНАЧЕНИЯ И СОКРАЩЕНИЯ
% ============================================================
\structsection{Обозначения и сокращения}

\begin{itemize}[leftmargin=1.25cm, labelsep=0.5cm, itemsep=0pt]
  \item[API] --- Application Programming Interface (программный интерфейс приложения)
  \item[REST] --- Representational State Transfer (передача состояния представления)
  \item[CRUD] --- Create, Read, Update, Delete (создание, чтение, обновление, удаление)
  \item[HTTP] --- HyperText Transfer Protocol (протокол передачи гипертекста)
  \item[JSON] --- JavaScript Object Notation (нотация объектов JavaScript)
  \item[SQL] --- Structured Query Language (язык структурированных запросов)
  \item[UUID] --- Universally Unique Identifier (универсальный уникальный идентификатор)
  \item[БД] --- база данных
  \item[ПО] --- программное обеспечение
  \item[СУБД] --- система управления базами данных
  \item[ОС] --- операционная система
  \item[DI] --- Dependency Injection (внедрение зависимостей)
  \item[FK] --- Foreign Key (внешний ключ)
  \item[PK] --- Primary Key (первичный ключ)
\end{itemize}

% ============================================================
% ВВЕДЕНИЕ
% ============================================================
\structsection{Введение}

В условиях современного информационного общества задачи планирования и организации деятельности приобретают всё большее значение. Эффективное управление задачами является ключевым фактором продуктивности как для отдельных пользователей, так и для команд. С ростом объёмов информации и количества параллельно выполняемых проектов возрастает потребность в специализированных программных средствах, позволяющих структурировать, отслеживать и анализировать ход выполнения задач.

Веб-приложения для управления задачами (Todo-приложения) представляют собой один из наиболее востребованных классов программного обеспечения. Существующие решения, такие как Todoist, Microsoft To~Do, TickTick и другие, предлагают широкий функционал, однако зачастую являются проприетарными, не предоставляют доступа к исходному коду и не позволяют адаптировать систему под специфические требования организации.

Актуальность настоящего исследования определяется необходимостью разработки веб-приложения для управления задачами с открытым исходным кодом, построенного на основе современных архитектурных подходов и технологий, обеспечивающих масштабируемость, поддерживаемость и расширяемость системы.

\textbf{Объект исследования}~--- веб-приложения для управления задачами.

\textbf{Предмет исследования}~--- методы и технологии проектирования RESTful веб-приложений с применением многослойной архитектуры на языке программирования Go.

\textbf{Цель работы}~--- исследование существующих методов и технологий для проектирования и разработки RESTful веб-приложения управления задачами с применением многослойной архитектуры.

Для достижения поставленной цели необходимо решить следующие \textbf{задачи}:

\begin{itemize}[leftmargin=1.25cm, itemsep=0pt]
  \item[---] провести анализ предметной области и обзор существующих аналогов;
  \item[---] исследовать архитектурные подходы к проектированию веб-приложений;
  \item[---] обосновать выбор языка программирования, СУБД и вспомогательных технологий;
  \item[---] исследовать применимые паттерны проектирования;
  \item[---] сформулировать техническое задание на разработку;
  \item[---] спроектировать структуру базы данных и REST~API;
  \item[---] определить тестовые сценарии для проверки стабильной работы системы;
  \item[---] подготовить входные данные для практической части.
\end{itemize}

\textbf{Методы исследования:} анализ научной и технической литературы, сравнительный анализ существующих решений, метод структурного проектирования, метод моделирования данных.

% ============================================================
% 1 АНАЛИЗ ПРЕДМЕТНОЙ ОБЛАСТИ
% ============================================================
\clearpage
\section{Анализ предметной области}

\subsection{Системы управления задачами: назначение и классификация}

Системы управления задачами (Task Management Systems) представляют собой класс программного обеспечения, предназначенного для планирования, отслеживания и контроля выполнения задач. Такие системы позволяют пользователям создавать задачи, назначать им приоритеты, устанавливать сроки выполнения, отмечать степень завершённости и получать статистику по продуктивности.

По масштабу применения системы управления задачами можно классифицировать следующим образом:

\begin{itemize}[leftmargin=1.25cm, itemsep=0pt]
  \item[---] персональные Todo-приложения, ориентированные на индивидуальное использование;
  \item[---] командные системы управления проектами, предоставляющие возможности совместной работы;
  \item[---] корпоративные системы, интегрированные с бизнес-процессами организации.
\end{itemize}

По архитектуре реализации выделяют: настольные приложения, веб-приложения, мобильные приложения и кроссплатформенные решения. В контексте настоящего исследования рассматриваются многопользовательские веб-приложения с REST~API интерфейсом.

\subsection{Обзор существующих аналогов}

Для определения требований к разрабатываемой системе проведён анализ существующих решений в области управления задачами.

\textbf{Todoist}~--- одно из наиболее популярных приложений для управления задачами с поддержкой более 30~миллионов пользователей. Предоставляет REST~API, поддерживает совместную работу и визуализацию продуктивности. Является проприетарным решением с закрытым исходным кодом.

\textbf{Microsoft To~Do}~--- бесплатное приложение от компании Microsoft, интегрированное с экосистемой Microsoft~365. Поддерживает синхронизацию между устройствами и REST~API через Microsoft~Graph.

\textbf{TickTick}~--- кроссплатформенное приложение с поддержкой календарного представления и таймера Pomodoro. Предоставляет REST~API.

\textbf{Google Tasks}~--- минималистичное приложение от Google, интегрированное с Gmail и Google~Calendar. Функционал ограничен базовыми операциями без статистики.

Результаты сравнительного анализа представлены в таблице~\ref{tab:analogs}.

\begin{table}[H]
\caption{Сравнительный анализ существующих аналогов}\label{tab:analogs}
\centering
\small
\begin{tabularx}{\textwidth}{|X|c|c|c|c|c|}
\hline
\textbf{Критерий} & \textbf{Todoist} & \textbf{MS To~Do} & \textbf{TickTick} & \textbf{Google Tasks} & \textbf{TodoApp} \\
\hline
Язык/Платформа & Python & C\#/.NET & Java & Java & Go \\
\hline
REST API & Да & Да & Да & Да & Да \\
\hline
Open Source & Нет & Нет & Нет & Нет & Да \\
\hline
Многопользоват. & Да & Да & Да & Да & Да \\
\hline
Статистика & Платно & Нет & Платно & Нет & Да \\
\hline
Swagger-док. & Нет & Нет & Нет & Нет & Да \\
\hline
Самост. развёрт. & Нет & Нет & Нет & Нет & Да \\
\hline
\end{tabularx}
\end{table}

\subsection{Выводы по анализу предметной области}

Проведённый анализ показал, что существующие решения обладают рядом существенных ограничений. Все рассмотренные аналоги являются проприетарными продуктами с закрытым исходным кодом. Функции статистики зачастую доступны только в платных версиях. Ни одно из решений не предоставляет автоматически генерируемую Swagger-документацию. Таким образом, обоснована необходимость разработки собственного веб-приложения с открытым исходным кодом.

% ============================================================
% 2 ИССЛЕДОВАНИЕ МЕТОДОВ И ТЕХНОЛОГИЙ
% ============================================================
\clearpage
\section{Исследование методов и технологий разработки}

\subsection{Архитектурные подходы к проектированию веб-приложений}

\subsubsection{Монолитная архитектура}

Монолитная архитектура представляет собой традиционный подход, при котором все компоненты приложения размещаются в единой кодовой базе и развёртываются как одно целое. По мере роста системы данный подход приводит к высокой связанности компонентов.

\subsubsection{Многослойная (layered) архитектура}

Многослойная архитектура предполагает разделение приложения на логические слои: слой представления (Presentation), слой бизнес-логики (Business Logic) и слой доступа к данным (Data Access). Данный подход обеспечивает разделение ответственности (Separation of Concerns).

\subsubsection{Чистая архитектура (Clean Architecture)}

Чистая архитектура, предложенная Р.~Мартином~\cite{martin2018}, развивает идеи многослойной архитектуры, вводя строгое правило направления зависимостей: зависимости всегда направлены от внешних слоёв к внутренним.

В разрабатываемом приложении применяется адаптированный вариант многослойной архитектуры с четырьмя слоями, представленный на рисунке~\ref{fig:architecture}.

\begin{figure}[H]
\centering
\begin{tikzpicture}[
  box/.style={rectangle, draw, minimum width=4cm, minimum height=1cm, text centered, font=\small},
  arr/.style={-{Stealth[length=3mm]}, thick}
]
\node[box] (transport) {Transport (HTTP)};
\node[box, below=0.8cm of transport] (service) {Service (бизнес-логика)};
\node[box, below=0.8cm of service] (repository) {Repository (доступ к данным)};
\node[box, below=0.8cm of repository] (database) {Database (PostgreSQL)};

\draw[arr] (transport) -- (service);
\draw[arr] (service) -- (repository);
\draw[arr] (repository) -- (database);
\end{tikzpicture}
\caption{Схема многослойной архитектуры приложения}\label{fig:architecture}
\end{figure}

\subsection{Обоснование выбора языка программирования Go}

Для реализации серверной части рассмотрены языки Go, Java, Python и Node.js. Результаты сравнения представлены в таблице~\ref{tab:languages}.

\begin{table}[H]
\caption{Сравнение языков программирования для разработки REST~API}\label{tab:languages}
\centering
\small
\begin{tabularx}{\textwidth}{|X|c|c|c|c|}
\hline
\textbf{Критерий} & \textbf{Go} & \textbf{Java} & \textbf{Python} & \textbf{Node.js} \\
\hline
Производительность & Высокая & Высокая & Средняя & Средняя \\
\hline
Потребление памяти & Низкое & Высокое & Среднее & Среднее \\
\hline
Конкурентность & Горутины & Потоки & GIL & Event loop \\
\hline
Стат. типизация & Да & Да & Нет & Нет \\
\hline
Станд. HTTP-сервер & Да & Нет & Нет & Да \\
\hline
Развёртывание & Бинарный файл & Требует JVM & Интерпретатор & Runtime \\
\hline
\end{tabularx}
\end{table}

Язык Go выбран по следующим основаниям~\cite{donovan2021, gospec}: встроенная поддержка конкурентности через горутины; компиляция в один статический бинарный файл; полноценный HTTP-сервер в стандартной библиотеке; строгая статическая типизация; низкое потребление памяти.

\subsection{Обоснование выбора СУБД PostgreSQL}

В качестве СУБД выбрана PostgreSQL версии~18~\cite{pgdocs}. Обоснование: открытый исходный код; поддержка UUID как типа данных; поддержка CHECK-ограничений; поддержка транзакций с уровнями изоляции; наличие высококачественного драйвера pgx для Go~\cite{pgx}.

\subsection{Технологии контейнеризации: Docker и Docker Compose}

Для обеспечения воспроизводимости среды применяются Docker и Docker~Compose~\cite{docker}. Docker~Compose используется для оркестрации сервисов: PostgreSQL, migrate, port-forwarder, swagger.

\subsection{Протокол REST и спецификация OpenAPI/Swagger}

REST~--- архитектурный стиль, предложенный Р.~Филдингом~\cite{fielding2000}. Спецификация OpenAPI~\cite{openapi} позволяет описать REST~API в машиночитаемом формате и автоматически генерировать документацию.

\subsection{Паттерны проектирования}

\subsubsection{Dependency Injection}

Внедрение зависимостей~--- паттерн, при котором зависимости передаются объекту извне. В приложении используется ручное внедрение: цепочка Repository~--- Service~--- Handler собирается в функции main.

\subsubsection{Repository Pattern}

Паттерн <<Репозиторий>>~\cite{fowler2020} инкапсулирует логику доступа к данным. Интерфейс репозитория определяется в пакете сервиса согласно принципу Interface Segregation~\cite{gamma1994}.

\subsubsection{Optimistic Locking}

Оптимистичная блокировка~--- при обновлении записи в условие WHERE включается текущая версия. При конфликте возвращается HTTP~409.

\subsubsection{Nullable-тип для PATCH-запросов}

Обобщённый тип Nullable[T] позволяет различать <<не передано>>, <<null>> и <<значение>> в PATCH-запросах.

% ============================================================
% 3 ТЕХНИЧЕСКОЕ ЗАДАНИЕ
% ============================================================
\clearpage
\section{Техническое задание на разработку}

\subsection{Общие сведения о системе}

Разрабатываемая система~--- многопользовательское веб-приложение TodoApp с REST~API, веб-интерфейсом и Swagger-документацией.

\subsection{Функциональные требования}

Функциональные требования представлены в таблице~\ref{tab:func_req}.

\begin{table}[H]
\caption{Функциональные требования к системе}\label{tab:func_req}
\centering
\small
\begin{tabularx}{\textwidth}{|c|X|c|}
\hline
\textbf{№} & \textbf{Требование} & \textbf{Приоритет} \\
\hline
1 & Создание пользователя (имя, телефон) & Высокий \\
\hline
2 & Получение списка пользователей с пагинацией & Высокий \\
\hline
3 & Получение пользователя по идентификатору & Высокий \\
\hline
4 & Обновление данных пользователя (PATCH) & Высокий \\
\hline
5 & Удаление пользователя & Высокий \\
\hline
6 & Создание задачи (заголовок, описание, автор) & Высокий \\
\hline
7 & Получение списка задач с фильтрацией & Высокий \\
\hline
8 & Обновление задачи (PATCH) & Высокий \\
\hline
9 & Удаление задачи & Высокий \\
\hline
10 & Автоматическая фиксация времени завершения & Средний \\
\hline
11 & Получение статистики по задачам & Средний \\
\hline
12 & Swagger-документация API & Средний \\
\hline
13 & Веб-интерфейс & Средний \\
\hline
\end{tabularx}
\end{table}

\subsection{Нефункциональные требования}

К системе предъявляются следующие нефункциональные требования: graceful shutdown с таймаутом 30~секунд; пул соединений с БД; валидация входных данных на транспортном и доменном уровнях; логирование запросов и ошибок; поддержка CORS; контейнеризация через Docker~Compose; управление схемой БД через миграции.

\subsection{Описание модулей системы}

Система разделена на модули, описанные в таблице~\ref{tab:modules}. Модульная диаграмма представлена на рисунке~\ref{fig:modules}.

\begin{table}[H]
\caption{Описание модулей системы}\label{tab:modules}
\centering
\small
\begin{tabularx}{\textwidth}{|l|X|}
\hline
\textbf{Модуль} & \textbf{Назначение} \\
\hline
core & Общие компоненты: домен, HTTP-сервер, middleware, логгер, пул БД \\
\hline
users & CRUD-операции с пользователями \\
\hline
tasks & CRUD-операции с задачами \\
\hline
statistics & Расчёт и получение статистики по задачам \\
\hline
web & Обслуживание веб-интерфейса \\
\hline
\end{tabularx}
\end{table}

\begin{figure}[H]
\centering
\begin{tikzpicture}[
  box/.style={rectangle, draw, minimum width=2.5cm, minimum height=0.8cm, text centered, font=\small},
  arr/.style={-{Stealth[length=2mm]}, thick}
]
\node[box] (main) {main.go};
\node[box, below left=1cm and 0.5cm of main] (users) {users};
\node[box, below=1cm of main] (tasks) {tasks};
\node[box, below right=1cm and 0.5cm of main] (stats) {statistics};
\node[box, right=0.3cm of stats] (web) {web};
\node[box, below=1.5cm of tasks] (core) {core};

\draw[arr] (main) -- (users);
\draw[arr] (main) -- (tasks);
\draw[arr] (main) -- (stats);
\draw[arr] (main) -- (web);
\draw[arr] (users) -- (core);
\draw[arr] (tasks) -- (core);
\draw[arr] (stats) -- (core);
\draw[arr] (web) -- (core);
\end{tikzpicture}
\caption{Модульная диаграмма системы}\label{fig:modules}
\end{figure}

\subsection{Проектирование базы данных}

База данных содержит две таблицы: users и tasks, связанные отношением <<один ко многим>>. ER-диаграмма представлена на рисунке~\ref{fig:er}.

\begin{figure}[H]
\centering
\begin{tikzpicture}[
  entity/.style={rectangle, draw, minimum width=4cm, minimum height=0.6cm, font=\small\ttfamily},
  title/.style={rectangle, draw, fill=gray!20, minimum width=4cm, minimum height=0.6cm, font=\small\bfseries\ttfamily}
]
\node[title] (ut) {users};
\node[entity, below=0pt of ut] (u1) {id : UUID (PK)};
\node[entity, below=0pt of u1] (u2) {version : INTEGER};
\node[entity, below=0pt of u2] (u3) {full\_name : VARCHAR(100)};
\node[entity, below=0pt of u3] (u4) {phone\_number : VARCHAR(15)};

\node[title, right=3cm of ut] (tt) {tasks};
\node[entity, below=0pt of tt] (t1) {id : UUID (PK)};
\node[entity, below=0pt of t1] (t2) {version : INTEGER};
\node[entity, below=0pt of t2] (t3) {title : VARCHAR(100)};
\node[entity, below=0pt of t3] (t4) {description : VARCHAR(1000)};
\node[entity, below=0pt of t4] (t5) {completed : BOOLEAN};
\node[entity, below=0pt of t5] (t6) {created\_at : TIMESTAMP};
\node[entity, below=0pt of t6] (t7) {completed\_at : TIMESTAMP};
\node[entity, below=0pt of t7] (t8) {author\_user\_id : UUID (FK)};

\draw[thick] (u1.east) -- ++(1cm,0) |- (t8.west) node[midway, above, font=\small] {1..*};
\end{tikzpicture}
\caption{ER-диаграмма базы данных}\label{fig:er}
\end{figure}

\subsection{Проектирование REST API}

Спроектированные эндпоинты представлены в таблице~\ref{tab:endpoints}.

\begin{table}[H]
\caption{Эндпоинты REST API}\label{tab:endpoints}
\centering
\small
\begin{tabularx}{\textwidth}{|c|l|X|l|}
\hline
\textbf{Метод} & \textbf{Путь} & \textbf{Описание} & \textbf{Коды} \\
\hline
GET & /api/v1/users & Список пользователей & 200, 500 \\
\hline
POST & /api/v1/users & Создать пользователя & 201, 400 \\
\hline
GET & /api/v1/users/\{id\} & Получить пользователя & 200, 404 \\
\hline
PATCH & /api/v1/users/\{id\} & Обновить пользователя & 200, 409 \\
\hline
DELETE & /api/v1/users/\{id\} & Удалить пользователя & 204, 404 \\
\hline
GET & /api/v1/tasks & Список задач & 200, 500 \\
\hline
POST & /api/v1/tasks & Создать задачу & 201, 400 \\
\hline
GET & /api/v1/tasks/\{id\} & Получить задачу & 200, 404 \\
\hline
PATCH & /api/v1/tasks/\{id\} & Обновить задачу & 200, 409 \\
\hline
DELETE & /api/v1/tasks/\{id\} & Удалить задачу & 204, 404 \\
\hline
GET & /api/v1/statistics & Статистика & 200, 500 \\
\hline
\end{tabularx}
\end{table}

% ============================================================
% 4 ТЕСТОВЫЕ СЦЕНАРИИ
% ============================================================
\clearpage
\section{Тестовые сценарии}

\subsection{Описание подхода к тестированию}

Для проверки стабильной работы определены тестовые сценарии, представленные в таблице~\ref{tab:tests}.

\begin{table}[H]
\caption{Тестовые сценарии}\label{tab:tests}
\centering
\small
\begin{tabularx}{\textwidth}{|c|X|X|X|}
\hline
\textbf{№} & \textbf{Сценарий} & \textbf{Входные данные} & \textbf{Ожидаемый результат} \\
\hline
1 & Создание пользователя & full\_name: <<Иван>>, phone: <<+79001234567>> & HTTP 201 \\
\hline
2 & Невалидный телефон & phone: <<12345>> & HTTP 400 \\
\hline
3 & Список пользователей & GET /api/v1/users & HTTP 200 \\
\hline
4 & Создание задачи & title, author\_user\_id & HTTP 201 \\
\hline
5 & Задача без автора & author\_user\_id: null & HTTP 400 \\
\hline
6 & Обновление статуса & PATCH completed: true & HTTP 200 \\
\hline
7 & Удаление задачи & DELETE /api/v1/tasks/\{id\} & HTTP 204 \\
\hline
8 & Несуществующая задача & GET /tasks/\{random\} & HTTP 404 \\
\hline
9 & Статистика & GET /api/v1/statistics & HTTP 200 \\
\hline
10 & Конфликт версии & PATCH с устаревшей версией & HTTP 409 \\
\hline
\end{tabularx}
\end{table}

% ============================================================
% 5 ВХОДНЫЕ ДАННЫЕ
% ============================================================
\section{Входные данные для практической части}

\subsection{Стек технологий}

Итоговый стек технологий представлен в таблице~\ref{tab:stack}.

\begin{table}[H]
\caption{Стек технологий}\label{tab:stack}
\centering
\small
\begin{tabularx}{\textwidth}{|X|l|l|X|}
\hline
\textbf{Компонент} & \textbf{Технология} & \textbf{Версия} & \textbf{Назначение} \\
\hline
Язык & Go & 1.24.1 & Основной язык \\
\hline
СУБД & PostgreSQL & 18.1 & Хранение данных \\
\hline
Контейнеризация & Docker & --- & Среда выполнения \\
\hline
Драйвер БД & pgx & 5.8.0 & Подключение к PostgreSQL \\
\hline
Логирование & zap & 1.27.1 & Структур. логирование \\
\hline
Swagger & swaggo/swag & 1.16.6 & Генерация OpenAPI \\
\hline
Валидация & validator & 10.30.1 & Валидация запросов \\
\hline
Миграции & golang-migrate & 4.19.1 & Управление схемой БД \\
\hline
\end{tabularx}
\end{table}

\subsection{Структура проекта}

Проект организован по принципу feature-based structure:

\begin{lstlisting}
cmd/todoapp/          - точка входа (main.go)
internal/core/        - общие компоненты ядра
internal/features/    - функциональные модули
  users/              - управление пользователями
  tasks/              - управление задачами
  statistics/         - статистика
  web/                - веб-интерфейс
migrations/           - SQL-миграции
docker-compose.yaml   - описание сервисов
Makefile              - команды управления
.env                  - переменные окружения
\end{lstlisting}

\subsection{Конфигурация окружения}

Конфигурация осуществляется через переменные окружения: HTTP\_ADDR~(:5050), HTTP\_SHUTDOWN\_TIMEOUT~(30s), POSTGRES\_USER, POSTGRES\_PASSWORD, POSTGRES\_DB, LOGGER\_LEVEL~(DEBUG), TIME\_ZONE~(UTC).

% ============================================================
% ЗАКЛЮЧЕНИЕ
% ============================================================
\structsection{Заключение}

В ходе выполнения научно-исследовательской работы получены следующие результаты.

Проведён анализ предметной области систем управления задачами. Рассмотрены аналоги: Todoist, Microsoft To~Do, TickTick, Google~Tasks. Выявлено, что все решения являются проприетарными.

Исследованы архитектурные подходы. Обоснован выбор многослойной архитектуры с элементами чистой архитектуры.

Обоснован стек технологий: Go, PostgreSQL, Docker, pgx, zap, swaggo, validator, golang-migrate.

Исследованы паттерны проектирования: Dependency Injection, Repository Pattern, Optimistic Locking, Nullable-тип.

Сформулировано техническое задание. Спроектирована структура БД и REST~API (11~эндпоинтов). Определены 10~тестовых сценариев. Подготовлены входные данные для практической реализации.

% ============================================================
% СПИСОК ИСПОЛЬЗОВАННЫХ ИСТОЧНИКОВ
% ============================================================
\structsection{Список использованных источников}

\begin{enumerate}[leftmargin=1.25cm, itemsep=2pt, label=\arabic*.]
  \item \label{src:donovan} Донован,~А. Язык программирования Go~/~А.~Донован, Б.~Керниган.~--- М.~: Вильямс, 2021.~--- 432~с.
  \item \label{src:martin} Мартин,~Р. Чистая архитектура. Искусство разработки программного обеспечения~/~Р.~Мартин.~--- СПб.~: Питер, 2018.~--- 352~с.
  \item \label{src:fielding} Fielding,~R.~T. Architectural Styles and the Design of Network-based Software Architectures~: дис. ...~д-ра философии~/~R.~T.~Fielding.~--- University of California, Irvine, 2000.~--- 162~p.
  \item \label{src:pgdocs} PostgreSQL~18 Documentation~[Электронный ресурс].~--- Режим доступа: https://www.postgresql.org/docs/18/ (дата обращения: 01.06.2026).
  \item \label{src:gospec} The Go Programming Language Specification~[Электронный ресурс].~--- Режим доступа: https://go.dev/ref/spec (дата обращения: 01.06.2026).
  \item \label{src:docker} Docker Documentation~[Электронный ресурс].~--- Режим доступа: https://docs.docker.com/ (дата обращения: 01.06.2026).
  \item \label{src:openapi} OpenAPI Specification~v3.0~[Электронный ресурс].~--- Режим доступа: https://swagger.io/specification/ (дата обращения: 01.06.2026).
  \item \label{src:gost732} ГОСТ~7.32-2017. Отчёт о научно-исследовательской работе.~--- М.~: Стандартинформ, 2017.
  \item \label{src:gost705} ГОСТ~Р~7.0.5-2008. Библиографическая ссылка.~--- М.~: Стандартинформ, 2008.
  \item \label{src:gamma} Gamma,~E. Design Patterns: Elements of Reusable Object-Oriented Software~/~E.~Gamma, R.~Helm, R.~Johnson, J.~Vlissides.~--- Addison-Wesley, 1994.~--- 395~p.
  \item \label{src:fowler} Фаулер,~М. Шаблоны корпоративных приложений~/~М.~Фаулер.~--- М.~: Вильямс, 2020.~--- 544~с.
  \item \label{src:richardson} Richardson,~L. RESTful Web Services~/~L.~Richardson, S.~Ruby.~--- O'Reilly Media, 2007.~--- 454~p.
  \item \label{src:pgx} pgx~--- PostgreSQL Driver and Toolkit for Go~[Электронный ресурс].~--- Режим доступа: https://github.com/jackc/pgx (дата обращения: 01.06.2026).
  \item \label{src:zap} Zap~--- Blazing fast, structured, leveled logging in Go~[Электронный ресурс].~--- Режим доступа: https://github.com/uber-go/zap (дата обращения: 01.06.2026).
  \item \label{src:migrate} golang-migrate~--- Database migrations written in Go~[Электронный ресурс].~--- Режим доступа: https://github.com/golang-migrate/migrate (дата обращения: 01.06.2026).
\end{enumerate}

% Для cite-ссылок в тексте
\begin{thebibliography}{99}
\bibitem{donovan2021} Донован~А., Керниган~Б. Язык программирования Go, 2021.
\bibitem{martin2018} Мартин~Р. Чистая архитектура, 2018.
\bibitem{fielding2000} Fielding~R.~T. Architectural Styles\ldots, 2000.
\bibitem{pgdocs} PostgreSQL~18 Documentation.
\bibitem{gospec} The Go Programming Language Specification.
\bibitem{docker} Docker Documentation.
\bibitem{openapi} OpenAPI Specification~v3.0.
\bibitem{gamma1994} Gamma~E. et~al. Design Patterns, 1994.
\bibitem{fowler2020} Фаулер~М. Шаблоны корпоративных приложений, 2020.
\bibitem{pgx} pgx --- PostgreSQL Driver for Go.
\end{thebibliography}

\end{document}
